%%%%%%%%%%%%%%%%%
% CV tailored for THALES - Ingénieur développement logiciel Segment Sol
% Strong positioning: Ground Software / Space Software Engineer | Python, Java, C, mission data, integration, IVVQ, CI/CD
%%%%%%%%%%%%%%%%%

\documentclass[8pt,a4paper,ragged2e,withhyper]{altacv}

\geometry{left=1.25cm,right=1.25cm,top=.5cm,bottom=1.cm,columnsep=1.2cm}

\usepackage{paracol}
\usepackage{fontawesome5}
\usepackage{hyperref}

\iftutex 
  \setmainfont{Roboto Slab}
  \setsansfont{Lato}
  \renewcommand{\familydefault}{\sfdefault}
\else
  \usepackage[rm]{roboto}
  \usepackage[defaultsans]{lato}
  \renewcommand{\familydefault}{\sfdefault}
\fi

\definecolor{SlateGrey}{HTML}{2E2E2E}
\definecolor{LightGrey}{HTML}{666666}
\definecolor{DarkPastelRed}{HTML}{8AC0D9}
\definecolor{PastelRed}{HTML}{00B0DA}
\definecolor{GoldenEarth}{HTML}{B4AD0C}

\colorlet{name}{black}
\colorlet{tagline}{PastelRed}
\colorlet{heading}{DarkPastelRed}
\colorlet{headingrule}{GoldenEarth}
\colorlet{subheading}{PastelRed}
\colorlet{accent}{PastelRed}
\colorlet{emphasis}{SlateGrey}
\colorlet{body}{LightGrey}

\definecolor{violet}{HTML}{5A2582}
\definecolor{rouge}{HTML}{F53B3B}
\definecolor{noir}{HTML}{00232C}
\definecolor{bleu}{HTML}{00B0DA}
\definecolor{rose}{HTML}{F831A0}
\definecolor{jaune}{HTML}{FFEC00}
\definecolor{vert}{HTML}{34AD6C}

\renewcommand{\namefont}{\Huge\rmfamily\bfseries}
\renewcommand{\personalinfofont}{\footnotesize}
\renewcommand{\cvsectionfont}{\LARGE\rmfamily\bfseries}
\renewcommand{\cvsubsectionfont}{\large\bfseries}
\renewcommand{\cvItemMarker}{{\small\textbullet}}
\renewcommand{\cvRatingMarker}{\faCircle}

\input{../../_shared/pubs-num.tex}

% auto_tex_manager layout tuning
\emergencystretch=3em
\hfuzz=60pt
\hbadness=9999
\sloppy

\begin{document}

\name{BOILEAU Guillaume}
\tagline{Ingénieur Logiciel Segment Sol | LISA, CNES/ESA, données L0--L3, cycle en V, Python \& IVVQ}

\photoL{2.8cm}{../../_shared/moi}

\personalinfo{%
  \email{guillaume.boileaupro@gmail.com}
  \phone{+33 6 59 94 85 84}
  \location{Alpes-Maritimes, FRANCE}
  \linkedin{guillaume-boileau-891b81265}
  \researchgate{Guillaume-Boileau-2}
  \orcid{0000-0002-3576-6968}
}

\makecvheader
\vspace{-0.3cm}

\AtBeginEnvironment{itemize}{\small}
\columnratio{0.64}

\begin{paracol}{2}

\cvsection{Experience}

\cvevent{\textbf{Software Engineer -- Développement logiciel, intégration \& support opérationnel}}{VEV Platform Services France}{2025 -- 2026}{Valbonne, France}

\textbf{Contexte :} Plateforme logicielle industrielle CPMS connectée à des infrastructures de recharge, avec services backend, interfaces applicatives, données opérationnelles, flux d’échanges et contraintes d’exploitation.

\textbf{Objectif :} Développer, intégrer et fiabiliser des services logiciels opérationnels, en intervenant sur les composants applicatifs, les interfaces de données, les incidents techniques, la validation des comportements et la documentation.

\begin{itemize}
  \item \textbf{Développement applicatif :} Contribution à des composants TypeScript, Angular et Node.js intégrés dans une plateforme opérationnelle.
  \item \textbf{Backend / frontend :} Travail sur des interactions entre interfaces utilisateur, services applicatifs, données opérationnelles et comportements métier.
  \item \textbf{Intégration de services :} Analyse des interfaces entre services logiciels, échanges de données, équipements connectés et contraintes terrain.
  \item \textbf{Flux opérationnels :} Investigation d’échanges FTP, interruptions de flux, incohérences de données et comportements inattendus.
  \item \textbf{Support technique :} Diagnostic d’incidents, analyse des causes, formalisation de constats et contribution aux actions de correction.
  \item \textbf{Validation opérationnelle :} Vérification de la cohérence entre données applicatives, comportement des équipements et attendus fonctionnels.
  \item \textbf{Pratiques projet :} Utilisation de Git, Linux, Zenhub, documentation technique et suivi collaboratif des sujets.
\end{itemize}

\divider

\cvevent{\textbf{Ingénieur R\&D senior -- Logiciel spatial, données mission \& validation}}{CNES}{2023 -- 2025}{Nice, France}

\textbf{Contexte :} Activités d’ingénierie logicielle et de R\&D pour le segment sol de la mission spatiale LISA ESA/NASA, impliquant Python scientifique, chaînes de données mission, niveaux L0--L3, intégration de modèles, simulation, validation technique et collaboration CNES/ESA.

\textbf{Objectif :} Concevoir et développer des briques logicielles de segment sol permettant d’intégrer, comparer, valider et exploiter des données et modèles hétérogènes dans le cadre du cycle en V et des chaînes L0--L3 de la mission LISA.

\begin{itemize}
  \item \textbf{Développement logiciel Python :} Conception et développement de CATpyLISA, plateforme dédiée à l’intégration, la comparaison, la validation et la revue de modèles.
  \textcolor{DarkPastelRed}{\href{https://gitlab.oca.eu/gboileau/lisa_l3_tools}{\bf GitLab: CATpyLISA}}
  \item \textbf{Segment sol LISA :} Développement de composants logiciels liés aux chaînes de données mission, à l’exploitation scientifique et à la validation des niveaux L0--L3.
  \item \textbf{Cycle en V \& niveaux de données :} Connaissance approfondie de la logique cycle en V appliquée aux chaînes de données spatiales, depuis les données brutes L0 jusqu’aux produits L3.
  \item \textbf{Intégration de modèles et produits :} Consolidation de contributions hétérogènes dans un cadre logiciel commun, traçable et reproductible, cohérent avec les besoins de validation du segment sol.
  \item \textbf{Validation / IVVQ :} Définition de contrôles de cohérence, règles de comparaison, critères d’acceptation et procédures de vérification sur des produits de données mission.
  \item \textbf{Analyse d’anomalies :} Identification d’écarts entre sorties, diagnostic des causes possibles, analyse d’impact et formalisation des points bloquants.
  \item \textbf{Documentation technique :} Rédaction de supports, synthèses, hypothèses, limites, résultats et éléments de revue collaborative.
  \item \textbf{Coordination technique :} Interface avec contributeurs scientifiques, logiciels et institutionnels dans un environnement spatial international.
\end{itemize}

\divider

\cvevent{\textbf{Data Scientist -- Python, traitement de données \& performance système}}{Universiteit Antwerpen, Belgium}{2021 -- 2023}{Antwerp, Belgium}

\textbf{Contexte :} Travaux de data science pour instruments de précision et détecteurs de nouvelle génération, combinant données capteurs, bruit environnemental, modélisation statistique, machine learning et validation de performance.

\textbf{Objectif :} Développer des workflows robustes pour analyser des données complexes, comparer des configurations, évaluer des performances et produire des résultats traçables.

\begin{itemize}
  \item \textbf{Pipelines Python :} Développement de workflows pour analyse de données, statistiques, performance, validation et reporting technique.
  \item \textbf{Données capteurs :} Analyse de mesures environnementales, propagation, corrélations spatiales et impacts sur la performance système.
  \item \textbf{Modélisation :} Mise en place de pipelines MCMC pour différentes configurations instrumentales et différents niveaux de bruit.
  \item \textbf{Machine learning :} Structuration d’approches de réduction de bruits non stationnaires avec canaux témoins, machine learning et deep learning.
  \item \textbf{Évaluation :} Utilisation de figures de mérite, analyses de sensibilité, contrôles de cohérence et comparaison de configurations.
  \item \textbf{Support décisionnel :} Production de résultats traçables, synthèses techniques et recommandations pour parties prenantes internationales.
\end{itemize}

\divider

\cvevent{\textbf{Ingénieur R\&D junior -- Calcul scientifique, simulation \& logiciel d’analyse}}{ARTEMIS, Observatoire de la Côte d’Azur}{2018 -- 2021}{Nice, France}

\textbf{Contexte :} Travaux de recherche appliquée liés à l’analyse de signaux faibles, à la simulation numérique, au traitement du signal et à la préparation de missions spatiales.

\textbf{Objectif :} Développer des méthodes logicielles d’analyse et de simulation afin de produire des résultats scientifiques robustes, documentés et exploitables.

\begin{itemize}
  \item \textbf{Calcul scientifique :} Développement de workflows d’analyse, simulation, visualisation et validation de résultats.
  \item \textbf{Simulation numérique :} Développement d’approches de simulation pour environnements complexes et grands jeux de données.
  \item \textbf{Traitement du signal :} Méthodes pour analyser et séparer des signaux faibles et superposés.
  \item \textbf{Validation :} Vérification de cohérence, comparaison de résultats, études de sensibilité et analyse de limites méthodologiques.
  \item \textbf{Documentation :} Formalisation des méthodes, hypothèses, résultats et conclusions pour revue collaborative.
\end{itemize}

\switchcolumn

\cvsection{Profile}

\textbf{Ingénieur docteur orienté développement logiciel segment sol, avec une expérience directe sur LISA, les chaînes de données mission L0--L3, le cycle en V, Python, l’intégration de modèles, l’IVVQ et la documentation technique en environnement CNES/ESA.}

Positionnement pour ce rôle : contribuer au développement de solutions logicielles sol pour programmes spatiaux, depuis l’analyse du besoin jusqu’à l’intégration, la validation, le support technique et l’amélioration continue.

\begin{itemize}
  \item Expérience directe du développement logiciel segment sol LISA, des chaînes de données mission, des niveaux L0--L3 et de la validation associée.
  \item Expérience industrielle en développement logiciel, frontend/backend, intégration, support opérationnel et analyse d’incidents.
  \item Culture spatiale forte : environnement CNES/ESA, cycle en V, données mission, modèles, simulation, validation, IVVQ et revues techniques.
  \item Capacité à collaborer avec équipes système, IVVQ, opérations, support, qualité, cybersécurité et produit.
  \item Bases solides en Git, Linux, Docker, CI/CD, documentation, tests, reproductibilité et anglais technique.
\end{itemize}

\cvsection{Languages}

\cvskill{English}{4}
\divider
\cvskill{Spanish}{2}
\divider

\medskip

\cvsection{Education}

\cvevent{PhD in Physics}{Université Côte d’Azur}{2018 -- 2021}{Nice, France}

\divider

\cvevent{Selective Graduate Program in Fundamental Physics (Magistère)}{Université Paris-Saclay}{2016 -- 2018}{Orsay, France}

\divider

\cvevent{MSc in Astronomy and Astrophysics}{Observatoire de Paris / Université Paris-Saclay}{2017 -- 2018}{Paris, France}

\divider

\cvevent{BSc in Fundamental Physics}{Université Paris-Saclay}{2015 -- 2016}{Orsay, France}

\cvsection{Professional Training}

\cvevent{Machine Learning in Python with scikit-learn}{FUN MOOC -- Inria}{2026}{France Université Numérique}

\small{
Predictive modelling, supervised learning, model evaluation, cross-validation, metrics, pipelines and practical machine learning workflows with Python and scikit-learn.
}

\divider

\cvevent{Supervised Machine Learning: Regression \& Classification}{DeepLearning.AI -- Andrew Ng}{2020}{Coursera}

\divider

\cvevent{Recherche reproductible : principes méthodologiques pour une science transparente}{FUN MOOC -- Inria}{2025}{France Université Numérique}

\divider

\cvevent{Continuous Professional Training -- Software, Data, AI and Systems}{OpenClassrooms}{2024 -- 2026}{}

\small{
Python, Java, SQL, Machine Learning, AI, Linux, JavaScript, TypeScript, Angular, Node.js, Django, REST APIs, C/C++, web development, algorithms and software engineering fundamentals.
}

\end{paracol}

\cvsection{Projects}

\cvevent{CATpyLISA / LISA Segment Sol -- Données mission L0--L3, intégration et validation}{Mission spatiale LISA ESA/NASA -- CNES/ESA}{2023 -- 2025}{}

Développement logiciel lié au segment sol de LISA : structuration de données mission, niveaux L0--L3, intégration de modèles, comparaison de résultats, validation technique et revue collaborative de workflows scientifiques.

\textbf{Rôle :} développement Python, structuration de données L0--L3, automatisation d’analyses, validation de modèles, analyse d’écarts, documentation technique et coordination avec plusieurs contributeurs.

\textbf{Compétences mobilisées :} Python, NumPy, SciPy, pandas, Matplotlib, GitLab, segment sol, cycle en V, données L0--L3, IVVQ, data quality, validation de modèles, simulation, documentation.

\divider

\cvevent{VEV -- Développement logiciel, intégration \& support opérationnel}{Industrial software / Backend-frontend / Data flows}{2025 -- 2026}{}

Développement et intégration de services logiciels pour une plateforme CPMS connectée à des infrastructures de recharge, avec analyse de flux FTP, investigation d’incidents, vérification de comportements applicatifs, support opérationnel et documentation technique.

\textbf{Rôle :} développement TypeScript, Angular et Node.js, intégration logicielle, analyse de flux FTP, investigation d’incidents, vérification de comportements applicatifs et formalisation de constats techniques.

\textbf{Compétences mobilisées :} TypeScript, Angular, Node.js, FTP, Linux, Git, analyse d’incidents, intégration logicielle, support opérationnel, documentation technique.

\divider

\cvevent{Workflows de simulation, traitement de données \& validation}{Scientific software / Mission preparation / Data processing}{2018 -- 2025}{}

Développement de workflows numériques pour traiter des données complexes, analyser des signaux faibles, comparer des modèles, évaluer des performances et produire des résultats traçables.

\textbf{Rôle :} calcul scientifique, simulation, traitement du signal, analyse de données, validation de résultats, diagnostic technique et reporting.

\textbf{Compétences mobilisées :} Python, C, signal processing, simulation, statistiques, MCMC, validation, analyse de performance, documentation technique.

\cvsection{Compétences clés}

\vspace{-.3cm}

\cvsubsection{Ground Software, spatial \& données mission}

{\small
\cvtag{Segment sol LISA}
\cvtag{Ground Software}
\cvtag{Systèmes spatiaux}
\cvtag{Données mission L0--L3}
\cvtag{Chaînes de traitement sol}
\cvtag{Supervision / monitoring}
\cvtag{Automatisation d’opérations -- awareness}
\cvtag{Centres de contrôle -- awareness}
\cvtag{Interfaces exploitation -- awareness}
\cvtag{LISA ESA/NASA}
\cvtag{CNES}
\cvtag{ESA}
\cvtag{Systèmes critiques}
\cvtag{Cycle en V}
\cvtag{Programmes spatiaux}
}

\divider\smallskip
\vspace{-.3cm}

\cvsubsection{Développement logiciel}

{\small
\cvtag{Python}
\cvtag{Java -- bases}
\cvtag{C}
\cvtag{TypeScript}
\cvtag{Angular}
\cvtag{Node.js}
\cvtag{Backend}
\cvtag{Frontend}
\cvtag{REST API}
\cvtag{SQL}
\cvtag{Scripts}
\cvtag{Applications métier}
\cvtag{Code review}
\cvtag{Documentation technique}
}

\divider\smallskip
\vspace{-.3cm}

\cvsubsection{Architecture, intégration \& cloud}

{\small
\cvtag{Microservices -- ramp-up}
\cvtag{Services distribués -- awareness}
\cvtag{Cloud -- ramp-up}
\cvtag{Docker}
\cvtag{Kubernetes -- ramp-up}
\cvtag{CI/CD}
\cvtag{Git}
\cvtag{GitLab}
\cvtag{Linux}
\cvtag{Bash}
\cvtag{Intégration logicielle}
\cvtag{Performance applicative -- awareness}
\cvtag{Scalabilité -- awareness}
\cvtag{Sécurité logicielle -- awareness}
}

\divider\smallskip
\vspace{-.3cm}

\cvsubsection{IVVQ, qualité \& support}

{\small
\cvtag{IVVQ}
\cvtag{Validation logicielle}
\cvtag{Validation système}
\cvtag{Tests fonctionnels}
\cvtag{Non-régression -- awareness}
\cvtag{Contrôles de cohérence}
\cvtag{Analyse d’anomalies}
\cvtag{Root cause analysis}
\cvtag{Support technique}
\cvtag{Support opérationnel}
\cvtag{Traçabilité}
\cvtag{Qualité logicielle}
\cvtag{Amélioration continue}
}

\divider\smallskip
\vspace{-.3cm}

\cvsubsection{Data, simulation \& analyse}

{\small
\cvtag{Data processing}
\cvtag{Simulation numérique}
\cvtag{Traitement du signal}
\cvtag{Analyse de performance}
\cvtag{Analyse de bruit}
\cvtag{Statistiques}
\cvtag{MCMC}
\cvtag{Machine Learning}
\cvtag{NumPy}
\cvtag{SciPy}
\cvtag{pandas}
\cvtag{Matplotlib}
\cvtag{Jupyter}
\cvtag{Data quality}
}

\divider\smallskip
\vspace{-.3cm}

\cvsubsection{Méthodes, collaboration \& outils}

{\small
\cvtag{Agile / Scrum}
\cvtag{Kanban}
\cvtag{Confluence}
\cvtag{Jira}
\cvtag{Zenhub}
\cvtag{OpenSearch}
\cvtag{Sentry}
\cvtag{Power BI}
\cvtag{Anglais technique}
\cvtag{Équipe pluridisciplinaire}
\cvtag{Interfaces système / IVVQ / opérations}
\cvtag{Curiosité technique}
\cvtag{Autonomie}
\cvtag{Rigueur}
}

\end{document}
